\chapter{Introduction}
\label{chap:introduction}

This dissertation was developed within the Doctoral Programme in Informatics Engineering at the Faculty of Engineering of the University of Porto. It investigates the development and operation of \gls{marl}-based control for distributed, heterogeneous, and dynamically reconfigurable \glspl{cps}, using smart-energy systems as the principal empirical domain.

This chapter introduces the context that motivates the research, establishes the transition from simulated learning to operational control, states the central thesis proposition and research questions, defines the scope of the work, summarises its main contributions, and presents the research path reflected in the organisation of the dissertation.

\section{Context and Motivation}
\label{sec:intro-context}

The research begins from a class of control problems in which decision-making is distributed across physical assets, software components, and participants whose capabilities and availability may change over time. The following sections move from this broader cyber-physical setting to Renewable Energy Communities and, finally, to the role of multi-agent reinforcement learning in decentralised control.

\subsection{Distributed and Dynamically Reconfigurable Cyber-Physical Systems}
\label{sec:intro-cps}

A \gls{cps} combines computation with physical processes. Sensors provide information about the physical environment, software interprets that information and decides what to do, and actuators change the physical process, creating a feedback loop between computation and the world. This connection is what distinguishes a \gls{cps} from a conventional software system: a software error is no longer confined to incorrect information on a screen; it can also lead to an incorrect physical action. Early work on \glspl{cps} already emphasised that physical dynamics, timing, safety, and reliability require abstractions that go beyond those used in general-purpose computing \parencite{lee2008cps}.

Many contemporary \glspl{cps} are distributed. A factory includes machines, local controllers, sensors, and supervisory services. A fleet-management system coordinates vehicles that move in and out of communication range. A group of robots may share a mission while retaining local sensing and control. An energy community connects buildings, meters, renewable generation, batteries, charging stations, and external platforms. In each case, information and decision authority are spread across several locations; no single controller has a complete view or exclusive ability to act.

Distribution changes the control problem in several ways. Each part of the system may observe only a local portion of the environment; communication introduces delays, bandwidth limits, and possible failures; available computation may range from cloud servers to embedded devices; and local decisions can have system-wide consequences. Local control can react quickly and preserve autonomy at the cost of a restricted view. Central control has access to broader context, with greater communication demands, scaling pressure, and exposure to a single point of failure.

The system may also be dynamically reconfigurable. In this thesis, \gls{dynamicreconfiguration} refers to changes in participating agents, available assets, capabilities, or interaction structure. Consider an electric vehicle connected to a household charger. Before the vehicle arrives, the charger has no energy request to satisfy. During charging, the vehicle's state, required departure charge, and available power become relevant to control. Once the vehicle leaves, that particular flexibility disappears even though the physical charger remains. Similar changes occur when a robot joins a mission, a machine is taken offline, a device fails, or a communication link becomes unavailable.

Such changes are different from ordinary variation in the value of a fixed observation. A temperature changing from \SI{20}{\celsius} to \SI{21}{\celsius} changes the state of an otherwise stable interface. A device or participant joining or leaving can instead change which observations exist, which actions are meaningful, which constraints apply, and which objectives can still be achieved. Control methods built around a fixed and homogeneous population may perform well in a stable benchmark yet require substantial adaptation when the configuration changes during operation.

Centralised optimisation remains a strong option when the system is sufficiently well characterised and the corresponding optimisation can be completed within the available decision time. Decentralised or hierarchical decision-making becomes increasingly relevant as configurations change, local information cannot be shared continuously, and communication or computation constrain central control. The central question is how control can remain coherent when information, capabilities, and decision authority are distributed across both local and shared mechanisms.

\subsection{Renewable Energy Communities as Representative Cyber-Physical Systems}
\label{sec:intro-recs}

The energy transition provides a concrete instance of these challenges. Electricity systems are incorporating increasing numbers of distributed and renewable resources, while electrification introduces flexible loads such as heat pumps and electric vehicles. Flexibility---the ability to shift production, consumption, storage, or exchange over time---is increasingly important for accommodating uncertainty and aligning demand with available renewable energy \parencite{hussain2023flexibility}.

European legislation defines a \gls{rec} as a legal entity based on open and voluntary participation, effectively controlled by members located near the renewable-energy projects it owns or develops, and primarily intended to provide environmental, economic, or social community benefits rather than financial profit \parencite{eu2018redii}. The control setting considered in this thesis is more concrete: a community may contain participants who consume and produce electricity and who expose different forms of flexibility through buildings, batteries, electric vehicles, chargers, heating systems, and other controllable loads.

A single participant may already combine assets with very different information and control interfaces. A battery can charge or discharge subject to power and \gls{soc} limits. A charger can exchange energy only while a compatible vehicle is connected and must account for the service expected at departure. A heat pump is coupled to thermal dynamics and comfort. These differences are not implementation details: they determine the observations, actions, constraints, and temporal dependencies of the control problem.

Time couples those decisions. Charging a battery now affects what it can do later. Delaying an electric-vehicle charge creates flexibility only until the vehicle must depart. Reducing heating demand in one interval may create a larger requirement later. Energy management is a sequential decision problem: the value of an action depends on its immediate effect and on the options it preserves for the future. This is one reason why \gls{rl} has attracted sustained interest in demand response and building control \parencite{vazquezcanteli2019demandresponse}.

At community level, objectives are also mixed. A participant may seek lower cost or guaranteed service, while the community may seek lower collective peaks, increased use of local renewable production, or adherence to a shared infrastructure limit. These objectives can align, but they are not identical. Several participants may benefit from coordination while simultaneously competing for a limited resource. This produces a mixed cooperative--competitive setting.

The same community can reconfigure during operation. Vehicles connect and disconnect, devices become unavailable, participants can join or leave, and external signals can change the value of different actions. Renewable Energy Communities bring distributed control, heterogeneous capabilities, changing availability, mixed objectives, and shared physical constraints into one empirical setting. These properties also arise in other \glspl{cps}, including robotics, transport, and industrial systems, although the empirical claims of this thesis remain grounded in the energy systems that are actually evaluated.

\begin{figure}[htbp]
\centering
\fbox{\parbox[c][4.4cm][c]{0.90\textwidth}{\centering
\textbf{Figure placeholder: Renewable Energy Community as a representative distributed CPS}\\[0.5em]
Participants and local agents; heterogeneous assets; local objectives; shared infrastructure; changing availability; community-level interactions.}}
\caption{Renewable Energy Community as the principal real-world cyber-physical-system instantiation used in the thesis.}
\label{fig:intro-rec-domain}
\end{figure}

\subsection{MARL for Mixed Cooperative--Competitive Decentralised Control}
\label{sec:intro-marl}

Rule-based and optimisation-based controllers remain important in energy management. A rule-based controller can encode transparent operational behaviour, while an optimisation model can schedule several assets from forecasts, explicit constraints, and models of their dynamics. These approaches provide strong references and, in some cases, appropriate operational fallbacks. Learning-based control becomes attractive when the decision strategy is difficult to specify completely in advance, when repeated interaction provides useful experience, or when decentralised decision makers must adapt to a changing environment.

\Gls{rl} formulates sequential control as a learning problem: an agent observes the environment, selects an action, receives feedback, and improves its policy through experience. The observations, actions, rewards, constraints, and training environment still have to be designed; learning does not remove the need for a sound control formulation. Its value lies in learning useful decision behaviour without encoding the complete strategy as a fixed rule set or solving a full optimisation problem at every decision step.

\Gls{marl} extends this setting to several agents interacting in a shared environment. It includes cooperative, competitive, and mixed formulations as well as different assumptions about information, communication, and training \parencite{albrecht2024marl,gronauer2022survey}. A common approach is \gls{ctde}, in which broader information can be used during training while each agent executes from its authorised local observations \parencite{lowe2017maddpg}. This is particularly relevant when operational control is distributed but learning can exploit richer information offline.

The mixed cooperative--competitive character is central to the energy setting. Two households with electric vehicles may share an interest in respecting a community capacity limit while each still prioritises the service required by its own vehicle. The agents must coordinate around shared outcomes without collapsing all participant objectives into one indistinguishable objective.

MARL introduces additional difficulties of its own. Other learning agents make the environment appear non-stationary from the perspective of an individual learner; joint decision spaces grow with the population; heterogeneous agents may expose different observations and actions; and communication or coordination can create additional computational and operational dependencies \parencite{hernandezleal2019survey,gronauer2022survey}. These difficulties become more pronounced when the physical system itself also changes. The learning algorithm is one part of the wider control system considered in this thesis, alongside its interfaces, physical constraints, runtime infrastructure, and evidence.

\section{Requirements for MARL in Real-World CPS}
\label{sec:intro-realworld-gap}

Simulation is indispensable for training and evaluating learning-based control. It enables repeated experiments, controlled comparisons, and the deliberate creation of disturbances that would be expensive, unsafe, or impractical to reproduce in a physical installation. At the same time, a simulator can provide cleaner observations, more stable configurations, and more computational resources than will exist during operation. Real-world use requires evidence that reaches beyond learning performance under those controlled conditions.

This thesis organises that transition around four cross-cutting requirements: \gls{scalability}, \gls{deployability}, \gls{safety}, and \gls{robustness}. These four requirements define the dimensions investigated explicitly across the architecture, algorithms, experimentation, lifecycle, and operational evidence developed in the thesis; they are not intended as a complete taxonomy of the challenges associated with real-world intelligent control.

\paragraph{Scalability.}
Scalability concerns more than increasing the number of agents. It includes the computational cost of learning and inference, growth in heterogeneous observations and actions, simulator throughput, communication and coordination overhead, the reproducible execution of many experiments, and the resources required by deployed components. A method that scales in one dimension can still fail in another; for example, parameter sharing can reduce model count while communication or critic cost continues to grow. Relevant evidence identifies the dimension that increased, the resources consumed, and the control quality maintained at that scale.

\paragraph{Deployability.}
A trained policy becomes an operational controller only when its experimental assumptions and runtime dependencies remain intact. It must stay associated with the simulator, scenario, interface, reward, and software versions under which it was produced; it must be validated for the intended target; and it must be packaged, installed, monitored, updated, and recoverable. In a distributed system, these concerns apply to compatible sets of runtime components, with no model file considered in isolation. Deployability in this thesis covers the controlled path from an experimental result to continued operation on heterogeneous edge and server-side infrastructure.

\paragraph{Safety.}
Physical control imposes limits that cannot be traded away for a higher reward. Batteries, chargers, buildings, and network connections have operating constraints, while users impose service requirements such as departure charge or comfort. Safety is used here in this operational sense: preventing unacceptable actions and states through the control formulation, validation, runtime constraints, and fallback mechanisms appropriate to each stage. Cybersecurity is important to real infrastructure but is not used as a synonym for safety in this thesis.

\paragraph{Robustness.}
Operational conditions do not remain identical to training conditions. Forecasts are imperfect, observations can be delayed or missing, devices can become unavailable, communication can fail, and the population or control interface can change. Robustness concerns whether useful control behaviour is maintained under these deviations and how performance degrades when assumptions no longer hold. Recovery after a disruption is treated as a related resilience concern and is addressed particularly through lifecycle and runtime mechanisms.

These requirements interact. More communication can improve coordination while increasing dependence on connectivity; a larger policy can improve control while exceeding an edge latency budget; conservative safety mechanisms can reduce violations while limiting performance; and frequent policy updates can support adaptation while increasing validation and operational risk. The thesis evaluates the four requirements jointly where their interactions matter and separately where causal attribution is needed.

The detailed state-of-the-art analysis in \cref{chap:foundations} examines how these challenges have been addressed in prior work and identifies the residual gaps that motivate the contributions developed in the remainder of the thesis. Other concerns are also important to real-world \glspl{cps}, including cybersecurity and privacy, human factors and interpretability, regulatory and market constraints, maintainability, and long-term organisational adoption. They are considered where they constrain the proposed methods or their operation, but they are not investigated directly as primary dimensions of this thesis.

\section{Thesis Proposition, Research Questions, and Scope}
\label{sec:intro-thesis-rqs-scope}

The research lies at the intersection of artificial intelligence, distributed systems, cyber-physical systems, edge computing, and intelligent energy control. Its central concern is computational: how distributed controllers can be represented, learned, evaluated, and operated when the physical and software interfaces through which they act are heterogeneous and can change. Energy is not treated as an interchangeable benchmark; asset dynamics, user requirements, shared infrastructure, and operational interfaces determine which decisions are meaningful and which evidence is required.

\subsection{Thesis Proposition}
\label{sec:intro-proposition}

The central proposition evaluated by this dissertation is:

\begin{quote}
\emph{Operationalising Multi-Agent Reinforcement Learning in distributed, heterogeneous, and dynamically reconfigurable cyber-physical systems requires the integrated co-design of suitable system representations and simulation environments, scalable learning and coordination mechanisms, safe and robust control processes, and lifecycle-aware distributed deployment infrastructure.}
\end{quote}

The proposition places the learning method within the system that makes its claims meaningful. A controller is trained against a representation of the physical system, evaluated under explicit assumptions, executed through concrete observation and action interfaces, and maintained through an operational lifecycle. The thesis investigates these elements together while evaluating the individual claims in the chapters where the corresponding evidence is produced.

\subsection{Research Questions}
\label{sec:intro-rqs}

Four research questions structure the investigation.

\paragraph{RQ1 -- Dynamic and Heterogeneous Multi-Agent Control.}
\begin{quote}
\emph{How can MARL-based control be formulated to accommodate heterogeneous and dynamically changing agent populations, capabilities, and availability, while balancing individual and collective objectives in distributed cyber-physical systems?}
\end{quote}
RQ1 concerns the representation and learning problem when agents and the entities they control are not assumed to expose one permanently fixed interface. It also addresses the coexistence of local and collective objectives under decentralised control.

\paragraph{RQ2 -- Multi-Dimensional Scalability of Learning and Coordination.}
\begin{quote}
\emph{Under which conditions, and through which algorithmic, computational, and communication mechanisms, can MARL-based learning and coordination maintain acceptable control performance and resource use as agent population size, heterogeneity, environment complexity, and dynamic reconfiguration increase, while computational and communication resources remain constrained?}
\end{quote}
RQ2 treats scale as a conditional, multi-dimensional property. It considers population size, heterogeneity, representation size, simulation and training cost, communication, inference, orchestration, and the resources available at deployment time.

\paragraph{RQ3 -- Safe and Robust Operation.}
\begin{quote}
\emph{How can physical and operational safety constraints, uncertainty, faults, dynamic reconfiguration, and shifts between training and operation be incorporated into and systematically assessed for MARL-based control?}
\end{quote}
RQ3 addresses both the prevention of unacceptable control actions and the maintenance of useful behaviour when observations, forecasts, devices, communication, or system configuration deviate from nominal conditions.

\paragraph{RQ4 -- Governed and Deployable MARL Policy-Set Lifecycle.}
\begin{quote}
\emph{What lifecycle abstractions and orchestration mechanisms allow coordinated MARL policy sets to be reproduced, governed, and deployed across heterogeneous and dynamically reconfigurable edge cyber-physical systems, while supporting compatibility validation, monitoring, staged updates, rollback, and failure recovery?}
\end{quote}
RQ4 follows the controller beyond training. It addresses experiment and environment lineage, compatibility between policies and runtime interfaces, validation and promotion, distributed deployment, monitoring, controlled updates, rollback, and recovery.

\subsection{Research Objectives}
\label{sec:intro-objectives}

The research questions are pursued through six objectives that connect the scientific argument to the experimental and systems work developed throughout the dissertation.

\begin{description}[style=nextline,leftmargin=1.2cm,labelindent=0cm]
  \item[O1 -- Characterise and formalise the control setting.]
  Establish the requirements of distributed, heterogeneous, and dynamically reconfigurable cyber-physical control, including mixed local and collective objectives and the four cross-cutting dimensions of scalability, deployability, safety, and robustness.

  \item[O2 -- Establish the simulation and experimental foundation.]
  Extend and validate a reproducible environment capable of representing heterogeneous energy assets, changing participants and capabilities, physical constraints, disturbances, and the scales required for controlled learning experiments.

  \item[O3 -- Develop and compare progressive control and learning solutions.]
  Design, implement, and evaluate rule-based, optimisation-based, reinforcement-learning, and multi-agent formulations under common scenarios and metrics, identifying what each approach solves and the assumptions under which it remains effective.

  \item[O4 -- Develop and evaluate TI-MARL.]
  Design a multi-agent method that operates over entity-bound typed interfaces whose observations, action capabilities, bounds, and health can change at runtime, and assess its behaviour under heterogeneity, reconfiguration, local feasibility constraints, faults, and scale.

  \item[O5 -- Develop the governed policy lifecycle and distributed runtime.]
  Define and evaluate the lineage, compatibility, validation, packaging, deployment, monitoring, update, rollback, and recovery mechanisms required to operate learned policies and policy sets across heterogeneous edge infrastructure.

  \item[O6 -- Integrate and validate the resulting approach in the energy domain.]
  Connect the architectural, simulation, control, learning, and lifecycle contributions within EnergAIze and evaluate their interactions using Renewable Energy Communities, electric-vehicle charging, real data, external platforms, edge execution, and physical demonstrators at the evidence levels actually available.
\end{description}


\subsection{Scope and Assumptions}
\label{sec:intro-scope}

The methods and system abstractions are developed for distributed and reconfigurable \glspl{cps}, while the empirical evidence is concentrated in smart-energy systems: Renewable Energy Communities, flexible buildings, electric vehicles, charging installations, and related energy services. Examples from robotics, transport, or industrial systems clarify the broader class of problems but do not constitute equivalent empirical validation. Generality concerns the formulated mechanisms and abstractions; measured performance remains specific to the evaluated energy configurations.

The thesis focuses on sequential operational control; long-term infrastructure planning and market design remain outside its central scope. Rule-based, optimisation-based, single-agent, and multi-agent controllers are all retained when they provide useful baselines, references, fallbacks, or solutions to particular subproblems. Simulation is used for training, controlled comparison, scaling, and disturbance injection, while real data, external platforms, edge devices, charging infrastructure, testbeds, and physical demonstrators provide additional evidence where available. The evidence classification keeps real-data simulation distinct from direct physical control.

The work assumes that observations can be obtained through an operational interface and that admissible actions can be applied to the controlled system. Lower-level device protections remain active. The thesis evaluates selected physical and operational constraints, disturbances, failures, and reconfiguration events, but it does not claim universal formal safety certification. Security, privacy, regulation, and human factors are considered where they affect the architecture or operation, but they are not primary research questions unless supported by dedicated methods and evidence.

\section{Research Contributions and Outcomes}
\label{sec:intro-contributions}

The thesis contributes scientific methods and findings together with the software and experimental infrastructure required to realise and evaluate them. The publications and research outputs associated with these contributions are listed in the front matter; detailed provenance is recorded in the chapters where collaborative work contributes directly to the corresponding evidence.

\subsection{Scientific Contributions}
\label{sec:intro-scientific-contributions}

The scientific contributions form seven related families. TI-MARL is the principal algorithmic contribution, while the remaining contributions establish the empirical and systems foundations, solve complementary control problems, and provide the evidence needed to evaluate the complete thesis proposition.

\paragraph{C1 -- Evidence-based characterisation of operational MARL requirements.}
The thesis develops a structured account of the conditions that separate conventional simulated MARL experiments from operational distributed \glspl{cps}. The literature analysis identifies existing solutions and counterexamples to simple absence claims, narrowing the residual challenge to the combined treatment of dynamic heterogeneity, multi-dimensional scalability, safe and robust operation, and a governed deployment lifecycle.

\paragraph{C2 -- The EnergAIze integrated research architecture.}
The thesis defines an architecture connecting data and scenario preparation, simulation, controller development, reproducible experimentation, policy governance, edge deployment, runtime monitoring, fallback, and recovery. Its contribution is the explicit organisation of responsibilities and contracts needed to study the path from a control problem to continued operation within one coherent research system.

\paragraph{C3 -- A physically audited simulation and experimentation environment.}
The CityLearn-derived environment is extended with heterogeneous assets, electric vehicles and chargers, electrical and community mechanisms, dynamic topology, entity-structured observations and actions, operational events, data and unit contracts, diagnostics, reproducible scenarios, and computational performance evidence. The contribution lies in their verified combination, in the class of experiments made possible by that combination, and in the continuity of evidence between training and deployment. The same versioned environment, scenario, interface, constraint, and evaluation contracts support experiment reconstruction, compatibility validation, policy promotion, packaging, and incident replay throughout the deployment lifecycle.

\paragraph{C4 -- Progressive control and learning solutions.}
The thesis develops and evaluates a progression of rule-based, optimisation-based, single-agent, and multi-agent control methods together with mechanisms for more stable, data-efficient, reusable, and hierarchical learning. Each stage addresses a concrete control problem and contributes empirical evidence about the assumptions under which it remains effective. These results provide both independent contributions and the design evidence from which the requirements of TI-MARL are derived.

\paragraph{C5 -- TI-MARL as the principal algorithmic contribution.}
TI-MARL is an optimizer-agnostic multi-agent learning method for systems in which the active agents, controlled entities, and observation--action interfaces can change. Versioned capability information and runtime health are compiled into entity-bound typed interfaces containing the observation parts, action-group instances, parameterised action ports, dependencies, validity conditions, and bounds currently available to each agent. A decentralised typed policy operates over this variable interface and produces one coherent local action bundle, while multi-agent coordination is learned through authorised shared observations, coupled outcomes, the joint physical transition, and variable-population central critics during training. The method is evaluated through controlled reconfiguration, health and disturbance tests, structural baselines, ablations, and multi-dimensional scaling.

\paragraph{C6 -- A governed lifecycle for MARL policies and policy sets.}
The thesis develops lifecycle abstractions linking experiments, environments, scenarios, interfaces, reward definitions, policies, compatible policy sets, validation states, packages, deployment targets, and runtime observations. The contribution focuses on lineage, compatibility-aware validation and promotion, distributed deployment, monitoring, staged updates, rollback, and recovery in multi-agent operational settings.

\paragraph{C7 -- Integrated real-world-oriented evidence in energy systems.}
The preceding contributions are connected through real data, external platforms, edge devices, charging infrastructure, energy communities, and fail-safe operation. The resulting evidence captures effects that do not appear in isolated learning experiments, including interface mismatches, runtime resource constraints, data-quality failures, deployment interventions, and the limits of the achieved operational evidence.

\subsection{System, Software, and Open-Source Contributions}
\label{sec:intro-artifacts}

The scientific work is supported by a set of interconnected research artefacts whose energy-domain integration is realised through \gls{energAIze}. These include data-collection and validation components; scenario and configuration tools; the extended simulation environment; rule-based, optimisation-based, single-agent, and multi-agent controllers; TI-MARL components; experiment workers and orchestration services; manifests and registries for experiments, policies, policy sets, packages, and deployment targets; validation and deployment tooling; edge-runtime components; monitoring, fallback, rollback, and recovery mechanisms; and reproducible configurations, documentation, and open-source releases where licensing permits.

These artefacts are not treated as substitutes for scientific evidence. They provide the implementation through which the corresponding claims can be tested, reproduced, integrated, and, where appropriate, transferred to operational infrastructure. Repository, release, licence, configuration, and experiment identifiers are retained for the submitted version, with access limitations stated where partner or data restrictions apply.

\subsection{Societal and Environmental Relevance}
\label{sec:intro-societal}

Renewable Energy Communities aim to give participants a more active role in the production, use, and sharing of energy. Intelligent control can help use local flexibility without requiring users to coordinate every flexible asset manually, while community-level objectives can support local renewable-energy use and reduce simultaneous peaks. The technical results of this thesis inform wider questions of flexibility, electrification, and distributed energy participation.

The thesis nevertheless distinguishes measured technical outcomes from broader societal claims. Lower simulated cost does not establish long-term adoption, peak reduction in a testbed does not establish system-wide emissions impact, and edge execution does not by itself establish a complete privacy guarantee. Societal and environmental implications are discussed from the evidence actually produced and within the limits of the evaluated systems.

\section{Research Path and Thesis Organisation}
\label{sec:intro-organisation}

The research developed cumulatively. Early work established the data, simulation, and experimental foundations and examined progressively more capable control formulations, from rule-based and optimisation approaches to single-agent and multi-agent reinforcement learning. Each stage addressed a concrete problem under explicit assumptions and produced results that remained useful in their own right. As the systems and experiments became more heterogeneous and dynamic, the remaining difficulty increasingly concerned how to preserve meaningful control when agents, assets, observations, actions, constraints, and operating conditions changed.

These findings motivated TI-MARL and, in parallel, the lifecycle and deployment mechanisms required to operate learned controllers beyond isolated experiments. The development was not a linear sequence: simulator limitations became visible during algorithm experiments, real operational data exposed assumptions in scenario design, and deployment targets introduced constraints that fed back into the control and software architecture. EnergAIze provides the common context in which these algorithmic and systems contributions are connected and evaluated.

The organisation of the dissertation follows this progression. \Cref{chap:foundations} establishes the foundations and residual gaps; \cref{chap:energAIze} presents the integrated architecture; and \cref{chap:simulator} develops the data, scenario, and simulation foundation. \Cref{chap:control-benchmark} presents the progression of control and learning solutions, leading to TI-MARL in \cref{chap:ti-marl}. \Cref{chap:lifecycle} addresses the governed lifecycle and distributed deployment of learned control, and \cref{chap:real-world} examines the integrated system in operational energy settings. Finally, \cref{chap:discussion} combines the evidence to answer the research questions, and \cref{chap:conclusions} concludes the thesis and identifies the remaining research directions.
